\subsection{函数的振幅}

与直径的概念相关的是函数 \(f\) （定义在任意集合 \(\mathbf{X}\) 上、取值于度量空间 \(\mathbf{X}'\) 中）的振幅的概念；若 \(\mathbf{A}\) 是 \(\mathbf{X}\) 的任意非空子集，则直径 \(\delta(f(\mathbf{A}))\) 称为 \(f\) 在 \(\mathbf{A}\) 中的振幅。

此外，若 \(\mathbf{X}\) 是拓扑空间 \(\mathbf{Y}\) 的子集，且 \(x \in \overline{\mathbf{X}}\) ，则数

\[
\omega (x; f) = \inf \delta (f (\mathbf {V} \cap \mathbf {X}))
\]

（当 V 取遍 \(x\) 在 Y 中的邻域滤子时）称为 \(f\) 在 \(x \in \overline{\mathbf{X}}\) 处的振幅。

命题 4. 振幅 \(\omega(x; f)\) ——任意函数 \(f\) （定义在拓扑空间 Y 的子集 X 上、取值于度量空间 \(\mathbf{X}'\) 中）的振幅—— 在 \(\overline{\mathbf{X}}\) 上是上半连续的。

设 \(a\) 是 \(\overline{\mathbf{X}}\) 中任意一点；则对每个 \(k > \omega(a; f)\) ，存在一个开邻域 \(\mathbf{V}\) （ \(a\) 的），使得 \(\delta(f(\mathbf{V} \cap \mathbf{X})) \leqslant k\) ；对每个 \(x \in \mathbf{V} \cap \overline{\mathbf{X}}\) ，\(\mathbf{V}\) 是 \(x\) 的一个邻域，因而

\[
\omega (x; f) \leqslant \delta (f (\mathbf {V} \cap \mathbf {X})) \leqslant k,
\]

这表明 \(\omega\) 在点 a 处上半连续。

为使 \(\omega(x; f) = 0\) 在点 \(x \in \overline{\mathbf{X}}\) 处成立，必要且充分的是：对每个 \(\varepsilon > 0\) ，存在一个邻域 \(\mathbf{V}\) （ \(x\) 的），使得 \(f(\mathbf{V} \cap \mathbf{X})\) 包含在半径为 \(\varepsilon\) 的一个球中；若 \(x \in \mathbf{X}\) ，这个条件表明 \(f\) 在点 \(x\) 处（关于 \(\mathbf{X}\) ）连续；若 \(x \in \overline{\mathbf{X}} \cap \complement \mathbf{X}\) ，则 \(f\) 映 \(\mathbf{X}\) 上的一个滤子基 —— \(x\) 在 \(\mathbf{Y}\) 中的邻域滤子的迹 —— 为 \(\mathbf{X}'\) 上的一个柯西滤子基；特别地：

命题 5. 设 \(f\) 是定义在一个子集 \(\mathbf{X}\) （拓扑空间 \(\mathbf{Y}\) 的子集）上、取值于完备度量空间 \(\mathbf{X}'\) 中的函数。则 \(f\) 关于 \(\mathbf{X}\) 在点 \(x \in \overline{\mathbf{X}}\) 处有极限的必要充分条件是 \(f\) 在 \(x\) 处的振幅为零。

